% !TEX root = ../my-thesis.tex

\chapter{Grundlagen wissenschaftlicher Workflows und Workflow-Management-Systeme}
\label{sec:basics}

\section{Wissenschaftliche Workflows}
\label{sec:basics:scientific-workflows}

\subsection{Begriff und Ziel wissenschaftlicher Workflows}
\label{sec:basics:scientific-workflows:definition}

\subsection{Typische Bestandteile eines Workflows}
\label{sec:basics:scientific-workflows:components}

\subsection{Wissenschaftliche Workflows in HPC-Umgebungen}
\label{sec:basics:scientific-workflows:hpc}

\section{Zentrale Konzepte wissenschaftlicher Workflows}
\label{sec:basics:concepts}

\subsection{Task}
\label{sec:basics:concepts:task}

\subsection{Abhängigkeiten und Directed Acyclic Graphs}
\label{sec:basics:concepts:dag}

\subsection{Scheduler und Batchsysteme}
\label{sec:basics:concepts:scheduler}

\subsection{Dataflow sowie file-based und stream-based Verarbeitung}
\label{sec:basics:concepts:dataflow}

\section{Typische Workflow-Strukturen und Muster}
\label{sec:basics:patterns}

\subsection{Sequence und Pipeline}
\label{sec:basics:patterns:pipeline}

\subsection{Parallel Split und Synchronization}
\label{sec:basics:patterns:parallel}

\subsection{Scatter-Gather/ Split-Apply-Combine}
\label{sec:basics:patterns:scatter}

\section{Workflow-Management-Systeme}
\label{sec:basics:wms}

\subsection{Aufgaben und Funktionsweise von Workflow-Management-Systemen}
\label{sec:basics:wms:tasks}

\subsection{Anforderungen an Workflow-Management-Systeme in HPC-Kontexten}
\label{sec:basics:wms:hpc-requirements}
\section{Die betrachteten Systeme}
\label{sec:basics:systeme}

\section{Verwandte Arbeiten}
\label{sec:basics:verwandte}

